\documentclass[dvipdfmx,a4paper]{article}
\usepackage{amsmath}
\usepackage{amssymb}
\usepackage{amsfonts}
\usepackage{bm}
\usepackage[utf8]{inputenc}
\usepackage{geometry}
\usepackage[dvipdfmx]{hyperref}
\usepackage{pxjahyper}
\geometry{left=25mm,right=25mm,top=25mm,bottom=25mm}

\title{なぜ電子・陽子・中性子だけが安定なのか：\\
二重時空におけるねじれ基準状態からの演繹的導出}

\author{https://github.com/hypernumbernet}
\date{\today}

\begin{document}

\maketitle

\begin{abstract}
二重時空理論（DST）において、すべての質量フェルミオンは、コンパクト双対時空のディオファントス量子化と二つの位相的不変量——ねじれ層の数と双対位相の巻き数——のもとでねじれ不整合スカラー \(J\) を最小化する双対ローター構成として実現される。本稿では、\(J\) の大域最小値とこれらの不変量が、三つの状態——電子（レプトン数 \(L=1\) を担う単層最小励起）、陽子（バリオン数 \(B=1\) を与える三層ねじれ星）、中性子（わずかに高いねじれエネルギーをもつ同一の三層構造）——によってのみ達成されることを示す。他のすべてのフェルミオン構成は、層数または巻き数の制約を破り厳密に正の \(\Delta J\) を生じて崩壊経路を開くか、双対角の離散性が課すディオファントス条件を満たさず、コンパクト双対幾何と両立しない。したがって観測される物質の安定性階層は、外部の保存則や背景連続体に頼ることなく、すべての質量粒子が運ぶ16次元双四元数構造の必然的な代数的帰結として従う。

\end{abstract}

\section{はじめに}

宇宙論的時間スケールで生き残るのが電子・陽子・（準安定な）中性子だけであり、他のすべての質量フェルミオンは崩壊するという経験的事実は、第一原理の水準では未だ説明されていない。標準模型では必要な保存則が手で課され、その微視的起源は未解決のままである。二重時空理論は、この欠けていた導出を与える。

双対時空がコンパクトであるため、双対ローター角は離散化される：
\[
\theta_a\in\frac{2\pi}{N}\mathbb{Z},\qquad N\propto\frac{m}{m_\text{Planck}}.
\]
ねじれ不整合スカラー \(J\) はしたがって離散的な最小値のみを許す。フェルミオン構成が安定であるための必要十分条件は、二つの位相的制約——(i) ねじれ層の数が整数であること、(ii) 双対位相の巻き数がその層数（またはその単純な倍数）と等しいこと——を満たしつつ \(J\) の絶対最小を実現することである。これらの条件を同時に満たすローター状態は、ちょうど三つだけである。

電子は、時間反転双対写像に関連するキラル射影と最小双対励起によって実現されるレプトン数 \(L=1\) の、ねじれ振動子の単層基準状態に対応する。そのねじれエネルギーは離散格子が許す絶対最小に位置する。

陽子は、双対位相巻き数3の三層ねじれ星を実現し、バリオン数 \(B=1\) と価数内容 \(uud\) を与える。三層共鳴はディオファントス境界を飽和し、すべての三層構成のうち \(J\) の大域最小を与える。

中性子は同一の三層位相トポロジーをもつが、ねじれエネルギーがわずかに高く、観測される質量差 \(m_n-m_p\approx1.293\,\mathrm{MeV}\) と整合する。この超過は、層構造自体が分単位の時間スケールで位相的に保護されたまま、双対セクターにおけるキラル弱結合が媒介する \(\beta\) 崩壊経路を開くのにちょうど必要な量である。

層数と巻き数の他のあらゆる組合せは、\(\Delta J>0\) を生じ（即時の崩壊経路を開く）か、離散量子化条件を破る（コンパクト双対幾何と両立しない状態となる）。したがって物質の安定性階層はパラメータ値の偶然ではなく、二重時空の内在的ねじれ動力学の一意の基準状態構造である。

本論文の残りの構成は次のとおりである。第2節では層数と巻き数によりすべての低エネルギーローター構成を分類し、上記三状態のみが \(J\) の大域最小を達成することを証明する。第3節では中性子のエネルギー超過を精密に導出し、\(\beta\) 崩壊が三層トポロジーを破壊せずに進行することを示す。第4節では実験的シグネチャと反証可能性を論じる。第5節で結論を述べる。

このように、物質安定性の長年の謎は純粋に代数的な言い換えに還元される：電子・陽子・中性子は、二重時空理論の離散双対格子上におけるねじれ基準状態問題の唯一の解である。

\section{低エネルギーローター構成の分類と大域的 J 最小性の証明}

二つの位相的整数——ねじれ層数 \(\ell \in \mathbb{N}\) と双対位相巻き数 \(w \in \mathbb{Z}\)——により、双対ローターのすべての低エネルギー構成を分類する。これらの整数は、離散量子化 \(\theta_a = 2\pi k/N\)（\(N \propto m/m_\text{Planck}\)）のもとでのコンパクト双対時空の既約表現をラベル付けする。ねじれ不整合スカラー \(J\) は、これらのラベルと連続的な不整合角 \(\delta\phi = \phi - \theta\) の関数である：
\[
J(\ell,w,\delta\phi) = \frac12 m \ell \,(\delta\phi)^2 + V_\text{top}(w-\ell),
\]
ここで位相的ポテンシャル \(V_\text{top}\) は \(w=\ell\)（または世代数 \(n\) をもつ複合状態では \(w = n\ell\)）のときゼロであり、それ以外では無限大である。ディオファントス条件はさらに、層共鳴条件
\[
D_N(\delta\phi) = \frac{\sin(N\delta\phi/2)}{N\sin(\delta\phi/2)}
\]
が有限かつ有界であるような許容 \(N\) に制限を課す。これはすべての整数 \(N\) で成立するが、\(J\) の離散的な最小を選び出す。

\(J\) の大域最小は、共鳴条件と双対写像 \(X\mapsto Xi\) が誘導するキラル割当てを満たしつつ、前因子 \(\ell\) を同時に最小化する \(\ell\) と \(w\) によってのみ達成される。

\begin{itemize}
\item \textbf{単層状態（\(\ell=1\)）}：許容される最小巻き数は \(w=1\) である。対応する状態は、時間反転双対セクターに関連するキラル射影 \(P_L = (1-i)/2\) によりレプトン数 \(L=1\) を担う。ねじれエネルギーは調和振動子の基準状態値に帰着する：
  \[
  J_{e} = \frac12 m_e \,(\delta\phi_{\rm min})^2,
  \]
  ここで \(\delta\phi_{\rm min}\) は格子が許す最小の不整合角である。この構成は電子によって実現される。\(w\neq1\) の他のすべての単層状態は \(V_\text{top}=\infty\) により禁止される。

\item \textbf{二層状態（\(\ell=2\)）}：共鳴と両立する最小巻き数は \(w=2\) である。しかし前因子 \(\ell=2\) は単層の場合に比べねじれ剛性を倍にし、
  \[
  J_{\ell=2} \ge 2\,J_e > J_e.
  \]
  を与える。この水準では安定なメソンやダイクォーク基準状態は存在せず、すべての二層構成は厳密に高い \(J\) に位置し、単層状態へ即座に崩壊する。

\item \textbf{三層状態（\(\ell=3\)）}：共鳴条件は \(w=3\) を唯一の許容巻き数として選ぶ。この構成はバリオン数 \(B=1\)（層あたりの巻き数が1に等しい）を担う。ねじれエネルギーは
  \[
  J_{\ell=3} = \frac32 m \,(\delta\phi_{\rm min})^2.
  \]
  である。このセクター内の基準状態は陽子（\(uud\)）であり、その価数内容は追加の不整合なしに三層共鳴を飽和する。同一位相トポロジーセクター内の第一励起状態は中性子（\(udd\)）であり、陽子との差は小さな同位スピン破れシフト \(\Delta J \approx 1.293\,\mathrm{MeV}/c^2\) のみである。この超過は、層トポロジーを保ったまま \(\beta\) 崩壊経路を開くのにちょうど必要な量である。

\item \textbf{高層状態（\(\ell\ge4\)）}：前因子 \(\ell\) は線形に増大するため、
  \[
  J_{\ell\ge4} \ge 4\,J_e \gg J_{\ell=3}.
  \]
  すべてのこのような状態は大域最小から遠く離れ、三層または単層基準状態への分裂に対して不安定である。
\end{itemize}

他の組合せがより低い \(J\) を与えないことを証明するため、層数 \(\ell'\) と巻き数 \(w'\) をもち \(J(\ell',w',\delta\phi') < J_{\rm min}\) となる構成が存在すると仮定する。すると、(i) \(w'\neq\ell'\)（世代因子を除く）であり、これは \(V_\text{top}=\infty\) を強制して除外されるか、(ii) \(w'=\ell'\) かつ \(\ell' < 3\) である。\(\ell'=1\) の場合、唯一の可能性は既に同定した電子である。\(\ell'=2\) では剛性の倍増が直ちに \(J> J_{\rm min}\) を与える。\(\ell'\ge4\) では前因子の線形増大が再び \(J_{\rm min}\) を超える。したがって大域最小は、上に列挙した三構成によってのみ達成される。

この分類は完備である：すべての質量フェルミオンは離散双対格子上の一意の組 \((\ell,w)\) に対応し、観測される安定性階層は、二重時空理論の位相的制約とディオファントス制約のもとでの \(J\) の最小化から直接従う。

\section{中性子のエネルギー超過とキラル弱結合による \texorpdfstring{$\beta$}{beta} 崩壊の代数的機構}

中性子と陽子は同一の三層トポロジー（\(\ell=3\)、\(w=3\)）を共有し、したがって同じバリオン数 \(B=1\) をもつ。両者の唯一の差は、双対ローターパラメータの小さな同位スピン破れシフトであり、これがねじれエネルギー超過
\[
\Delta J = J_n - J_p \approx 1.293\,\mathrm{MeV}/c^2.
\]
を生じる。この超過は、価数内容 \(udd\)（中性子）が、双対写像が誘導するキラル割当てを保ちつつ離散双対格子上で共鳴条件を満たすために、\(uud\)（陽子）よりわずかに大きな双対位相不整合 \(\delta\phi_n\) を要求するために生じる。シフトはパラメトリックに小さく、電磁細構造定数とQCDスケールの積のオーダーであるが、運動学的に許される崩壊経路を開くのにちょうど必要な量である。

二重時空理論において弱相互作用は、双対セクターにおけるキラルゲージ結合として実現される。完全双対ローターは最小結合
\[
R_\text{dual} \to \mathcal{P}\exp\Bigl(\frac12\int(\theta_a + g_W W_a P_L)\Gamma_a\,ds\Bigr),
\]
により拡張される。ここで \(P_L = (1-i)/2\) は時間反転双対写像 \(X\mapsto Xi\) によって生成されるキラル射影である。左手性射影は \(udd\to uud\) 遷移のみを選び出し、右手性成分は結合しない。したがって崩壊演算子は同位スピン射影を変化させるが、層数 \(\ell=3\) と巻き数 \(w=3\) は不変のままである。

遷移振幅はしたがって
\[
\mathcal{M}(n\to p+e^-\bar\nu_e) \propto g_W \langle p | P_L | n \rangle \exp(-S_\text{tors}),
\]
ここでねじれ作用 \(S_\text{tors}\) は、二つの三層状態を結ぶ最小 \(J\) 経路に沿って評価される。経路が同一位相セクター内に留まるため、指数抑制は穏やか（\(\exp(-\Delta J/\hbar\omega)\)）のみであり、観測される寿命 \(\tau_n \approx 880\,\mathrm{s}\) を与える。層数または巻き数を変化させるいかなる崩壊も、無限の位相的障壁 \(V_\text{top}\) により厳密に禁止される。

このように中性子は準安定である：三層トポロジーは保護されるが、離散双対格子が許す小さなねじれ超過により、陽子基準状態への遅いキラル媒介遷移が可能となる。この機構は完全に代数的であり、外部の弱ボゾン場や即興の選択則は不要である。標準模型においてパリティ破れを強制するのと同じキラル構造が、ここでは双対時空の時間反転性質の直接的幾何的帰結として現れる。

\section{実験的反証可能性と具体的予測}

本枠組みは、標準模型および他の量子重力アプローチの両方と区別する、鋭い反証可能な予測の集合を与える。すべての予測は、ねじれ層の位相的保護と双対時空の離散性から直接従い、追加パラメータは導入されない。

\subsection{異種安定フェルミオンの不在}

層数 \(\ell \neq 1,3\) または巻き数 \(w \neq \ell\) をもついかなるフェルミオン構成も、無限の位相的障壁 \(V_\text{top}\) により禁止されるか、\(\Delta J \gg m_e c^2\) に位置する。したがって本理論は、安定または長寿命の異種フェルミオン（例：安定テトラクォーク、ペンタクォーク、既知三世代を超える重いレプトン）の完全な不在を予測する。LHC または将来の加速器でこのような状態が観測されれば、本枠組みは直ちに反証される。

\subsection{中性子 \texorpdfstring{$\beta$}{beta} 崩壊観測量の精密測定}

キラル双対結合は、中性子崩壊における電子--反中微子角相関係数 \(a\) とファイアツ干渉項 \(b\) に、小さいが計算可能な補正を導入する。崩壊が保護された三層セクター内で完全に進行するため、補正は \(\Delta J / m_n c^2 \approx 1.4 \times 10^{-3}\) のオーダーである。PERKEO III、UCNA、または今後の PERKEO IV および Nab 実験による高統計測定は、今後2年以内にこの精度に到達しうる。\(10^{-4}\) 水準での null 結果は、中性子--陽子質量差のねじれ的起源を排除する。

\subsection{ねじれ星からの重力波エコー}

三層（またはそれ以上）のねじれ構成から成るコンパクト天体——「ねじれ星」——は、引力層と斥力層が交互に重なる層状内部をもつ。このような天体に入射する重力波は、各ねじれ界面で部分的に反射し、特徴的なエコー列を生じる。連続するエコー間の時間遅延は
\[
\Delta t \sim \frac{GM}{c^3} \approx 10^{-5}\Bigl(\frac{M}{10\,M_\odot}\Bigr)\,\mathrm{s},
\]
であり、位相変調は双対ロータースペクトルによって固定される。LISA と Einstein Telescope は、連星合体または銀河中心からのこのエコーに感度をもつ。5年間のデータ取得後も非観測であれば、強重力の離散的ねじれ構造は反証される。

\subsection{プランク期キラル重力波}

双対時間矢のプランク期固定（第1節）は、テンソル--スカラー比 \(r \sim \delta^2 \approx 10^{-10}\) とパリティ奇なスペクトルをもつキラル確率的重力波背景を源とする。この背景は、特徴的な TB および EB 相関により LiteBIRD または CMB-S4 によって検出可能である。予測振幅での非零キラル信号の検出は、バリオン非対称性の幾何的起源、ひいては電子・陽子・中性子を安定化するねじれ機構の直接的証拠となる。

四つの予測はすべて、離散双対格子スケールが観測される中性子--陽子質量差によって固定されればパラメータフリーである。したがってこれらは、物質の安定性階層が二重時空理論の必然的代数的帰結であるという主張の決定的検証を構成する。

\section{結論}

我々は、物質の経験的安定性階層——電子と陽子の絶対的安定性および中性子の準安定性——が、二重時空理論の必然的代数的帰結として従うことを示した。すべての質量フェルミオンは、コンパクト双対時空の離散格子上の双対ローター構成によって特徴づけられる。ねじれ不整合スカラー \(J\) は、二つの位相的不変量——層数 \(\ell\) と双対位相巻き数 \(w\)——のもとで最小化される。\(J\) の大域最小は、次の三構成によってのみ達成される：

\begin{itemize}
\item 電子（\(\ell=1\)、\(w=1\)、\(L=1\)）、
\item 陽子（\(\ell=3\)、\(w=3\)、\(B=1\)）、
\item 中性子（小さなねじれ超過 \(\Delta J \approx 1.293\,\mathrm{MeV}/c^2\) をもつ同一の三層トポロジー）。
\end{itemize}

他のすべての組合せは、無限の位相的障壁 \(V_\text{top}\) により禁止されるか、厳密に高い \(J\) に位置して即時の崩壊経路を開く。バリオン数およびレプトン数保存は、物理的に許されるいかなる遷移においても層数と巻き数が保存されるという要求として直接現れる。それらはもはや人為的に課されるのではなく、16次元双四元数代数の幾何の帰結である。

中性子の長いが有限の寿命は、小さなエネルギー超過が保護された三層セクター内での遅いキラル媒介遷移を許す一方、トポロジーを変化させるいかなる過程も厳密に禁止されるために生じる。この機構は完全に代数的であり、離散双対格子スケールが観測される中性子--陽子質量差によって固定されればパラメータフリーである。

本枠組みは四つの鋭い反証可能な予測を与える：異種安定フェルミオンの完全な不在、\(10^{-3}\) 水準での中性子 \(\beta\) 崩壊観測量への計算可能な補正、ねじれ星からの重力波エコー、およびプランク期からのキラル確率的重力波背景。これらの予測は、今後10年以内に次世代実験によって検証される。

要約すれば、なぜ電子・陽子・中性子だけが唯一の長寿命質量粒子なのかという長年の問いに答えが与えられた：それらは二重時空の内在的ねじれ動力学の一意の基準状態である。一般相対性理論と標準模型はそれぞれの領域で回復されるが、物質の安定性は経験的偶然ではなく幾何的必然として明らかになる。連続体仮説は放棄される。最も深いレベルでの実在は離散的で粒子局所的かつ二重時間的である。

\end{document}
